\chapter{Spaces, Quotients and Duality}\label{ch:duality}

\begin{center}\itshape
I have every confidence that the effort I have applied to the science reported upon here, which has occupied a
considerable span of my lifetime and demanded the most intense exertions of my powers, is not to be lost.
\upshape(Hermann Grassmann, foreword to the \emph{Ausdehnungslehre} of 1862)
\end{center}

In 1844 a schoolteacher from Stettin published a book that nobody could read. Hermann Grassmann had never studied
mathematics at a university; he had studied theology and philology, taught at a secondary school, and worked out on his
own a theory of ``extensive magnitudes'' that could be added and stretched, in any number of dimensions. His
\emph{Ausdehnungslehre} contained, in philosophical prose, vector spaces, linear independence, dimension, subspaces and
their sums and intersections, the exterior product, and the formula $\dim(U+W)+\dim(U\cap W)=\dim U+\dim W$. It sold
almost no copies. M\"obius, asked to review it, confessed he could not get through it. Grassmann applied for a
university chair and was turned down; the referee, Kummer, praised the content and found the form impossible. In 1862 he rewrote the whole book in a more mathematical style. It sold even fewer copies. Grassmann
then turned to Sanskrit, translated the Rig Veda, and wrote a dictionary of it that is still in print. The linguists
gave him an honorary doctorate; the mathematicians caught up only after his death, when Peano, in 1888, extracted from
Grassmann's work the axioms of a vector space that every student now learns.

This chapter revisits the beginning of linear algebra in the spirit of Grassmann: not matrices first, but spaces and
the maps between them. We rebuild bases and dimension, then meet two constructions that a first course often skips and
that the rest of the book uses constantly. The \emph{quotient} $V/W$ makes precise the idea of ``computing modulo $W$''.
The \emph{dual space} $V^*$ is the space of all linear measurements of $V$, and it turns every statement about a map into
a statement about its \emph{transpose}. The chapter ends with a method that is pure Grassmann: to bound the size of a
family of objects, turn them into vectors and show that they are linearly independent. Their number can then not
exceed the dimension.

\begin{think}
A town has $32$ inhabitants who love forming clubs, but the town council imposes two rules: every club must have an
\emph{odd} number of members, and any two different clubs must have an \emph{even} number of members in common. How many
clubs can there be at most? Can you find $32$? Can you find $33$? Section~\ref{sec:dimarg} settles the question in
four lines.
\end{think}

\section{Spaces and bases}

Throughout, $\mathbb K$ is a field: $\Q$, $\R$, $\C$ or the finite field $\mathbb F_p=\Z/p\Z$, and later any finite field.
A \emph{vector space} over $\mathbb K$ is a set $V$ with an addition making it an abelian group and a multiplication by
scalars satisfying $1\cdot v=v$, $\lambda(\mu v)=(\lambda\mu)v$, $(\lambda+\mu)v=\lambda v+\mu v$ and
$\lambda(v+w)=\lambda v+\lambda w$. These are Peano's axioms, and everything in this book follows from them. The
examples to keep in mind are $\mathbb K^n$, the matrices $M_{m,n}(\mathbb K)$, the polynomials $\mathbb K[X]$ and
$\mathbb K[X]_{\le n}$ (degree at most $n$), and the space $\mathbb K^S$ of all functions from a set $S$ to $\mathbb K$.

A family $(v_i)_{i\in I}$ is \emph{free} (linearly independent) if the only finite linear combination
$\sum\lambda_iv_i$ that vanishes is the one with all $\lambda_i=0$; it \emph{spans} $V$ if every vector is a finite
linear combination of the $v_i$; it is a \emph{basis} if both hold, that is, if every vector is in exactly one way a
finite combination of the $v_i$. The whole theory of dimension rests on one lemma.

\begin{lemma}[Steinitz exchange]\label{lem:steinitz}
If $v_1,\dots,v_m$ are linearly independent in the span of $w_1,\dots,w_n$, then $m\le n$, and after renumbering the
$w_j$, the family $v_1,\dots,v_m,w_{m+1},\dots,w_n$ spans the same space as $w_1,\dots,w_n$.
\end{lemma}
\begin{proof}
By induction on $m$, the case $m=0$ being empty. Assume the statement for $v_1,\dots,v_{m-1}$: so $m-1\le n$ and
$v_1,\dots,v_{m-1},w_m,\dots,w_n$ spans the space $E=\operatorname{span}(w_1,\dots,w_n)$. Write
$v_m=\sum_{i<m}\alpha_iv_i+\sum_{j\ge m}\beta_jw_j$. If all $\beta_j$ were zero (in particular if $m-1=n$ and there is no
$w_j$ left), $v_m$ would be a combination of $v_1,\dots,v_{m-1}$, contradicting independence. So $m\le n$ and some
$\beta_j\ne0$; renumber so that $\beta_m\ne0$. Then $w_m=\beta_m^{-1}\bigl(v_m-\sum_{i<m}\alpha_iv_i-\sum_{j>m}\beta_jw_j\bigr)$
lies in the span of $v_1,\dots,v_m,w_{m+1},\dots,w_n$, which therefore contains all of $v_1,\dots,v_{m-1},w_m,\dots,w_n$
and so equals $E$.
\end{proof}

\begin{theorem}[dimension]\label{thm:dim}
Let $V$ be spanned by finitely many vectors. Then $V$ has a finite basis, and all bases have the same number of elements,
the \emph{dimension} $\dim V$. Every free family can be completed into a basis, every spanning family contains a basis,
and in a space of dimension $n$ a family of $n$ vectors is free if and only if it spans.
\end{theorem}
\begin{proof}
Take a spanning family $w_1,\dots,w_n$ and remove vectors that are combinations of the others, one at a time, until none
is: the remaining family spans and is free (a vanishing combination with some $\lambda_i\ne0$ would express that $v_i$
through the others). Two bases $B,B'$ have $|B|\le|B'|$ and $|B'|\le|B|$ by Lemma~\ref{lem:steinitz}. A free family
$v_1,\dots,v_m$ completes into a basis by the exchange lemma applied with a basis $w_1,\dots,w_n$: the family
$v_1,\dots,v_m,w_{m+1},\dots,w_n$ spans and has $n$ elements, and a spanning family of $n=\dim V$ elements is free,
because a spanning family contains a basis, which has $n$ elements, so it is the whole family. The same remark shows that
a free family of $n$ vectors, completed into a basis of $n$ vectors, was already a basis.
\end{proof}

Grassmann's formula follows at once: take a basis of $U\cap W$, complete it into a basis of $U$ and, separately, into a
basis of $W$; the union of the three pieces is a basis of $U+W$ (check it!), which gives
\[
\dim(U+W)+\dim(U\cap W)=\dim U+\dim W.
\]

\begin{remark}
A space that is not spanned by finitely many vectors is \emph{infinite-dimensional}. Every vector space still has a
basis, but the proof needs Zorn's lemma (Tome~6, Chapter~1), and the basis cannot in general be written down: nobody
has ever exhibited a basis of $\R$ as a vector space over $\Q$. That $\R$ is infinite-dimensional over $\Q$ is easy,
though: a $\Q$-space of finite dimension $n$ is in bijection with $\Q^n$, hence countable, while $\R$ is not. In this
book infinite dimension appears only in such remarks; its proper home is Tome~7.
\end{remark}

\section{Sums, complements and quotients}\label{sec:quotients}

Subspaces $W_1,\dots,W_k$ of $V$ form a \emph{direct sum} $W_1\oplus\dots\oplus W_k$ if every vector of $W_1+\dots+W_k$ is
in exactly one way a sum $w_1+\dots+w_k$ with $w_i\in W_i$; for two subspaces this means $W_1\cap W_2=\{0\}$. A
\emph{complement} of $W$ in $V$ is a subspace $C$ with $V=W\oplus C$. In finite dimension complements exist: complete a
basis of $W$ into a basis of $V$ and let $C$ be the span of the added vectors. But there are many complements (every
line through the origin other than $W$ is a complement of a line $W$ in the plane), and none is better than the others.
The quotient is the canonical object that all complements are copies of.

\begin{definition}\label{def:quotient}
Let $W$ be a subspace of $V$. Two vectors are \emph{congruent modulo $W$}, $v\equiv v'\pmod W$, if $v-v'\in W$. This is
an equivalence relation; the class of $v$ is the affine subspace $v+W=\{v+w:w\in W\}$. The set of classes is the
\emph{quotient} $V/W$, with the operations $(v+W)+(v'+W)=(v+v')+W$ and $\lambda(v+W)=\lambda v+W$.
\end{definition}

The operations are well defined: if $v\equiv u$ and $v'\equiv u'$ then $(v+v')-(u+u')=(v-u)+(v'-u')\in W$ and
$\lambda v-\lambda u\in W$. The axioms of a vector space hold in $V/W$ because they hold in $V$, and the map
$\pi:V\to V/W$, $v\mapsto v+W$, is linear and onto, with kernel $W$. Figure~\ref{fig:quotient} shows the picture to keep
in mind: when $W$ is a line in the plane, the points of $V/W$ are the lines parallel to $W$, and they are added and
scaled like the points where they cross any fixed complement.

\begin{figure}[ht]
\centering
\fig{ch01-quotient}
\caption{The quotient of the plane by a line $W$: its elements are the lines parallel to $W$. Any complement $C$ meets
each of them once, so $C\to V/W$ is an isomorphism; the quotient itself does not depend on the choice of $C$.}
\label{fig:quotient}
\end{figure}

\begin{theorem}[quotients]\label{thm:quotient}
Let $W\subseteq V$, with $V$ finite-dimensional.
\begin{enumerate}[label=(\alph*)]
\item $\dim V/W=\dim V-\dim W$; more precisely, for every complement $C$ of $W$, the restriction $\pi|_C:C\to V/W$ is an
isomorphism.
\item (Universal property.) A linear map $f:V\to U$ with $W\subseteq\ker f$ factors in a unique way as $f=\bar f\circ\pi$
with $\bar f:V/W\to U$ linear.
\item (First isomorphism theorem.) For every linear $f:V\to U$, $\bar f:V/\ker f\to\operatorname{Im}f$ is an
isomorphism. In particular $\dim V=\dim\ker f+\operatorname{rank}f$.
\item (Third isomorphism theorem.) If $W\subseteq W'\subseteq V$, then $W'/W$ is a subspace of $V/W$ and
$(V/W)/(W'/W)\cong V/W'$.
\end{enumerate}
\end{theorem}
\begin{proof}
(a) If $c\in C$ and $\pi(c)=0$ then $c\in W\cap C=\{0\}$, so $\pi|_C$ is injective. Every $v$ writes $w+c$, and
$\pi(v)=\pi(c)$, so $\pi|_C$ is onto. Hence $\dim V/W=\dim C=\dim V-\dim W$.

(b) The only possible definition is $\bar f(v+W)=f(v)$, and it is well defined: if $v-v'\in W$ then $f(v)-f(v')=f(v-v')=0$.
It is linear because $f$ is.

(c) By (b) with $W=\ker f$, $\bar f$ exists and has image $\operatorname{Im}f$. It is injective: $\bar f(v+\ker f)=0$ means
$f(v)=0$, that is, $v+\ker f$ is the zero class. The dimension formula is (a).

(d) $W'/W=\pi(W')$ is a subspace. The map $V/W\to V/W'$, $v+W\mapsto v+W'$, is well defined by (b) (applied to the
projection $V\to V/W'$, whose kernel $W'$ contains $W$), onto, with kernel $\{v+W:v\in W'\}=W'/W$; apply (c).
\end{proof}

The quotient is how mathematicians ``compute modulo'' something. Three examples will recur.

\begin{example}\label{ex:quotients}
(1) The congruences of Tome~1 are quotients of groups; the vector space version is $\mathbb K[X]/(P)$, the
polynomials modulo the multiples of a polynomial $P$ of degree $d$. Euclidean division shows that every class contains
exactly one polynomial of degree $<d$, so $\mathbb K[X]_{<d}$ is a complement of $(P)$ and $\dim\mathbb K[X]/(P)=d$.
For $P=X^2+1$ over $\R$, the class of $X$ squares to $-1$: this quotient is $\C$, and this is how Cauchy defined the
complex numbers in 1847.

(2) Let $S$ be a set, $T\subseteq S$ and $W=\{f\in\mathbb K^S:f|_T=0\}$. The restriction map $\mathbb K^S\to\mathbb K^T$
is onto with kernel $W$, so $\mathbb K^S/W\cong\mathbb K^T$: a function modulo those vanishing on $T$ is just its
restriction to $T$.

(3) A linear system $Ax=b$ with $A\in M_{m,n}(\mathbb K)$ is solvable if and only if $b\in\operatorname{Im}A$, that is,
if and only if $b$ is zero in the \emph{cokernel} $\mathbb K^m/\operatorname{Im}A$. The dimension of the cokernel,
$m-\operatorname{rank}A$, counts the independent conditions on $b$. Section~\ref{sec:transpose} will tell us which
conditions.
\end{example}

\section{The dual space}

\begin{definition}\label{def:dual}
A \emph{linear form} on $V$ is a linear map $V\to\mathbb K$. The linear forms make a vector space $V^*$, the \emph{dual}
of $V$. If $e_1,\dots,e_n$ is a basis of $V$, the \emph{coordinate forms} $e_1^*,\dots,e_n^*$ are defined by
$e_i^*(e_j)=\delta_{ij}$: so $e_i^*(v)$ is the $i$-th coordinate of $v$.
\end{definition}

A linear form is a linear \emph{measurement}: a way to extract one number from a vector, compatibly with addition and
scaling. Evaluating a polynomial at a point, integrating a function, taking the trace of a matrix, reading the $i$-th
coordinate: all are linear forms.

\begin{proposition}[dual basis]\label{prop:dualbasis}
If $e_1,\dots,e_n$ is a basis of $V$, then $e_1^*,\dots,e_n^*$ is a basis of $V^*$, the \emph{dual basis}, and every
$\varphi\in V^*$ equals $\sum_i\varphi(e_i)\,e_i^*$. In particular $\dim V^*=\dim V$.
\end{proposition}
\begin{proof}
Both sides of $\varphi=\sum_i\varphi(e_i)e_i^*$ are linear and take the value $\varphi(e_j)$ at each $e_j$, so they are
equal; hence the $e_i^*$ span. If $\sum_i\lambda_ie_i^*=0$, evaluating at $e_j$ gives $\lambda_j=0$.
\end{proof}

So $V$ and $V^*$ are isomorphic, but the isomorphism $e_i\mapsto e_i^*$ depends on the basis: change $e_1$ into $2e_1$
and $e_1^*$ becomes $\tfrac12e_1^*$. There is no natural way to turn a vector into a measurement (without extra
structure, such as the inner product of Chapter~4). Remarkably, the other direction is natural: a vector \emph{is} a
measurement of measurements.

\begin{theorem}[the double dual]\label{thm:bidual}
For finite-dimensional $V$, the map $\iota:V\to V^{**}$, $\iota(v)(\varphi)=\varphi(v)$, is an isomorphism, and it is
defined without choosing a basis.
\end{theorem}
\begin{proof}
$\iota(v)$ is linear in $\varphi$, and $\iota$ is linear in $v$. If $\iota(v)=0$, then $\varphi(v)=0$ for all
$\varphi$; completing $v\ne0$ into a basis $v=e_1,e_2,\dots$ we would get $e_1^*(v)=1$, so $v=0$. Thus $\iota$ is
injective between spaces of the same dimension $n$ (Proposition~\ref{prop:dualbasis} twice), hence bijective.
\end{proof}

\begin{remark}
In infinite dimension, $\iota$ is still injective but never onto, and $V^*$ is strictly bigger than $V$: for instance the
dual of $\mathbb K[X]$ is the space $\mathbb K^{\N}$ of all sequences (a form is determined by its values on
$1,X,X^2,\dots$, which are arbitrary), and for $\mathbb K=\Q$ the first is countable and the second is not. Functional
analysis (Tome~7) restricts to \emph{continuous} linear forms to recover a good theory.
\end{remark}

Dual bases have a computational meaning: \emph{to find the coordinates of a vector in a basis, find the dual basis.}
Here is the most famous instance.

\begin{example}[Lagrange interpolation]\label{ex:lagrange}
Let $x_0,\dots,x_n\in\mathbb K$ be distinct and $V=\mathbb K[X]_{\le n}$. The evaluations $\varepsilon_i(P)=P(x_i)$ are
$n+1$ linear forms on a space of dimension $n+1$. The \emph{Lagrange polynomials}
\[
L_j(X)=\prod_{k\ne j}\frac{X-x_k}{x_j-x_k}
\]
satisfy $\varepsilon_i(L_j)=\delta_{ij}$. Hence the $L_j$ are linearly independent (apply $\varepsilon_i$ to a vanishing
combination), so they form a basis of $V$, and $(\varepsilon_i)$ is its dual basis. Proposition~\ref{prop:dualbasis}
read in $V$ gives the interpolation formula: every $P$ of degree at most $n$ is $P=\sum_jP(x_j)L_j$. In particular,
values $y_0,\dots,y_n$ can be prescribed arbitrarily at the $x_i$ by a unique polynomial of degree at most $n$.
\end{example}

The dual of a matrix space is described by the trace. For $A\in M_{n,m}(\mathbb K)$, let $\tau_A(X)=\operatorname{tr}(AX)$
for $X\in M_{m,n}(\mathbb K)$.

\begin{proposition}[the trace pairing]\label{prop:trace}
The map $A\mapsto\tau_A$ is an isomorphism $M_{n,m}(\mathbb K)\to M_{m,n}(\mathbb K)^*$: every linear form on matrices is
$X\mapsto\operatorname{tr}(AX)$ for a unique $A$.
\end{proposition}
\begin{proof}
$\operatorname{tr}(AX)=\sum_{i,j}a_{ij}x_{ji}$, so $\tau_A(E_{ji})=a_{ij}$ for the elementary matrices $E_{ji}$. Hence
$\tau_A=0$ forces $A=0$, and the map is an injective linear map between spaces of the same dimension $mn$.
\end{proof}

\section{Annihilators and the transpose}\label{sec:transpose}

\begin{definition}
The \emph{annihilator} of a subset $S\subseteq V$ is $S^\circ=\{\varphi\in V^*:\varphi(s)=0\text{ for all }s\in S\}$, the
measurements that do not see $S$. For a linear map $f:V\to U$, the \emph{transpose} $f^{\mathsf T}:U^*\to V^*$ is
$f^{\mathsf T}(\psi)=\psi\circ f$: to measure $v$ through $f$, measure $f(v)$.
\end{definition}

\begin{theorem}[annihilators]\label{thm:annihilator}
Let $V$ have dimension $n$ and $W\subseteq V$ be a subspace.
\begin{enumerate}[label=(\alph*)]
\item $\dim W^\circ=n-\dim W$, and $W^\circ\cong(V/W)^*$ canonically.
\item Under the identification $V=V^{**}$ of Theorem~\ref{thm:bidual}, $W^{\circ\circ}=W$. Equivalently, $W$ is the set of
common zeros of the forms vanishing on $W$: every subspace of codimension $k$ is an intersection of $k$ hyperplanes
$\ker\varphi_1\cap\dots\cap\ker\varphi_k$.
\item $(W_1+W_2)^\circ=W_1^\circ\cap W_2^\circ$ and $(W_1\cap W_2)^\circ=W_1^\circ+W_2^\circ$.
\end{enumerate}
\end{theorem}
\begin{proof}
(a) A form vanishing on $W$ factors through $\pi:V\to V/W$ (Theorem~\ref{thm:quotient}(b)), and conversely each
$\bar\varphi\in(V/W)^*$ gives $\bar\varphi\circ\pi\in W^\circ$; these maps are inverse to each other. So
$\dim W^\circ=\dim(V/W)^*=n-\dim W$.

(b) Clearly $W\subseteq W^{\circ\circ}$ (each $w$ kills every form killing $W$), and by (a) applied twice
$\dim W^{\circ\circ}=n-(n-\dim W)=\dim W$. For the second statement, take a basis $\varphi_1,\dots,\varphi_k$ of
$W^\circ$ ($k=n-\dim W$): a vector killed by all $\varphi_i$ is killed by all of $W^\circ$, so lies in $W^{\circ\circ}=W$.

(c) The first equality is immediate. For the second, $W_1^\circ+W_2^\circ\subseteq(W_1\cap W_2)^\circ$, and both have
dimension $n-\dim(W_1\cap W_2)$, by Grassmann's formula in $V^*$, by (a), and by the first equality:
$\dim(W_1^\circ+W_2^\circ)=(n-\dim W_1)+(n-\dim W_2)-(n-\dim(W_1+W_2))=n-\dim(W_1\cap W_2)$.
\end{proof}

\begin{theorem}[the transpose]\label{thm:transpose}
Let $f:V\to U$ be linear between finite-dimensional spaces.
\begin{enumerate}[label=(\alph*)]
\item If $f$ has matrix $A$ in bases $(e_j)$ of $V$ and $(u_i)$ of $U$, then $f^{\mathsf T}$ has matrix $A^{\mathsf T}$ in
the dual bases.
\item $\ker f^{\mathsf T}=(\operatorname{Im}f)^\circ$ and $\operatorname{Im}f^{\mathsf T}=(\ker f)^\circ$.
\item $\operatorname{rank}f^{\mathsf T}=\operatorname{rank}f$. In matrix terms: the row rank of a matrix equals its column
rank.
\end{enumerate}
\end{theorem}
\begin{proof}
(a) The $(j,i)$ entry of the matrix of $f^{\mathsf T}$ is the $e_j$-coordinate of $f^{\mathsf T}(u_i^*)$, namely
$f^{\mathsf T}(u_i^*)(e_j)=u_i^*(f(e_j))=a_{ij}$.

(b) $\psi\in\ker f^{\mathsf T}$ means $\psi(f(v))=0$ for all $v$, that is, $\psi\in(\operatorname{Im}f)^\circ$. Next, if
$\varphi=\psi\circ f$ then $\varphi$ kills $\ker f$, so $\operatorname{Im}f^{\mathsf T}\subseteq(\ker f)^\circ$. By the
first part and Theorem~\ref{thm:annihilator}(a),
$\operatorname{rank}f^{\mathsf T}=\dim U-\dim(\operatorname{Im}f)^\circ=\operatorname{rank}f$, while
$\dim(\ker f)^\circ=\dim V-\dim\ker f=\operatorname{rank}f$. Equal dimensions give equality, and this proves (c) too.
\end{proof}

The equality of row rank and column rank, which looks like a coincidence of computation when proved by Gaussian
elimination, is here a structural fact: the transpose is the same map seen from the side of measurements. Part (b) is the
precise answer to Example~\ref{ex:quotients}(3).

\begin{corollary}[Fredholm alternative]\label{cor:fredholm}
For $A\in M_{m,n}(\mathbb K)$ and $b\in\mathbb K^m$, the system $Ax=b$ has a solution if and only if $y^{\mathsf T}b=0$ for
every $y$ with $A^{\mathsf T}y=0$. Exactly one of the following happens: $Ax=b$ is solvable, or there is $y$ with
$A^{\mathsf T}y=0$ and $y^{\mathsf T}b=1$.
\end{corollary}
\begin{proof}
Identify $(\mathbb K^m)^*$ with row vectors $y^{\mathsf T}$. By Theorem~\ref{thm:transpose}(b),
$\ker A^{\mathsf T}=(\operatorname{Im}A)^\circ$, and by Theorem~\ref{thm:annihilator}(b), $\operatorname{Im}A$ is the set
of vectors killed by $(\operatorname{Im}A)^\circ$. If some $y\in\ker A^{\mathsf T}$ has $y^{\mathsf T}b\ne0$, rescale it so that
$y^{\mathsf T}b=1$.
\end{proof}

The certificate $y$ is worth noticing. If a system has a solution, exhibiting it convinces anyone; the alternative says
that if it has none, there is also a short convincing proof of that: a combination of the equations that reads $0=1$.
This is the seed of duality in linear programming (Tome~C2) and of Delsarte's bounds in Chapter~9.

\begin{example}[Lights Out]\label{ex:lightsout}
In the electronic game \emph{Lights Out} (Tiger Electronics, 1995), a $5\times5$ grid of lights is given, some on and
some off. Pressing a light toggles it and its horizontal and vertical neighbours (Figure~\ref{fig:lightsout}). The goal
is to switch all lights off. Over $\mathbb F_2$, pressing a set of buttons $x\in\mathbb F_2^{25}$ toggles the
configuration $Ax$, where $A$ is the symmetric $25\times25$ matrix with $a_{ij}=1$ when $i=j$ or $i,j$ are neighbours.
Pressing twice is like not pressing, and the order does not matter, so a configuration $b$ can be switched off if and only
if $Ax=b$ is solvable. A computation (by hand with patience, or in the notebook) shows that $\ker A$ has dimension $2$,
spanned by two ``quiet patterns'' $q_1,q_2$ that change nothing when pressed. By the Fredholm alternative and
$A^{\mathsf T}=A$:
\[
b\text{ is solvable}\iff q_1\cdot b=q_2\cdot b=0\quad\text{in }\mathbb F_2 .
\]
So exactly a quarter of the $2^{25}$ configurations can be solved, a parity check on $b$ tells which, and each solvable
configuration has exactly $4$ solutions, differing by quiet patterns. On a $3\times3$ or $6\times6$ board, $A$ is
invertible and every configuration is solvable; on $4\times4$ the kernel has dimension $4$.
\end{example}

\begin{figure}[ht]
\centering
\fig{ch01-lightsout}
\caption{Lights Out. Left: pressing the marked button toggles it and its four neighbours. Right: one of the two quiet
patterns of the $5\times5$ board; pressing all its lit buttons changes nothing, and every solvable configuration has an
even number of lit lights in common with it.}
\label{fig:lightsout}
\end{figure}

\section{The dimension argument}\label{sec:dimarg}

We come to the method announced at the beginning: \emph{to bound the size of a family, find a space in which its
members are linearly independent.} It is one of the most powerful ideas in combinatorics, and Tome~C1, Chapter~11
develops it at length. Here are three instances, the first of which answers the opening question.

\begin{theorem}[Oddtown; Berlekamp, 1969]\label{thm:oddtown}
Let $C_1,\dots,C_m$ be subsets of $\{1,\dots,n\}$ such that each $|C_i|$ is odd and each $|C_i\cap C_j|$, $i\ne j$, is
even. Then $m\le n$.
\end{theorem}
\begin{proof}
Let $v_i\in\mathbb F_2^n$ be the indicator vector of $C_i$. With the dot product $u\cdot v=\sum u_kv_k$ computed in
$\mathbb F_2$, $v_i\cdot v_j=|C_i\cap C_j|\bmod2$, which is $1$ if $i=j$ and $0$ otherwise. If $\sum_i\lambda_iv_i=0$,
taking the dot product with $v_j$ gives $\lambda_j=0$. So the $v_i$ are independent in $\mathbb F_2^n$ and $m\le n$.
\end{proof}

So the town of $32$ can have $32$ clubs (the $32$ single-member clubs will do) but never $33$. The proof is the same
as the one that showed the Lagrange polynomials independent: a family $(v_i)$ together with measurements $(\varphi_j)$
such that $\varphi_j(v_i)=\delta_{ij}$. It works just as well when the matrix $(\varphi_j(v_i))$ is only triangular with
non-zero diagonal; this \emph{triangular criterion} is the workhorse of the method.

\begin{proposition}[triangular criterion]\label{prop:triangular}
Let $v_1,\dots,v_m\in V$ and $\varphi_1,\dots,\varphi_m\in V^*$ satisfy $\varphi_i(v_i)\ne0$ and $\varphi_i(v_j)=0$ for $j<i$.
Then $v_1,\dots,v_m$ are linearly independent, and so are $\varphi_1,\dots,\varphi_m$.
\end{proposition}
\begin{proof}
Suppose $\sum_j\lambda_jv_j=0$ with some $\lambda_j\ne0$, and let $k$ be the largest index with $\lambda_k\ne0$. Apply
$\varphi_k$: the terms $j<k$ vanish by hypothesis and the terms $j>k$ have $\lambda_j=0$, leaving
$\lambda_k\varphi_k(v_k)=0$, a contradiction. For the forms, suppose $\sum_i\mu_i\varphi_i=0$ with some $\mu_i\ne0$, let $k$
be the smallest index with $\mu_k\ne0$ and evaluate at $v_k$: the terms $i>k$ vanish and the terms $i<k$ have $\mu_i=0$.
\end{proof}

The second instance is a cleverly chosen dual basis. A \emph{hyperplane} of $V$ is a subspace of codimension $1$, that is,
the kernel of a non-zero linear form (Theorem~\ref{thm:annihilator}(b) with $k=1$).

\begin{theorem}\label{thm:hyperplane}
For $n\ge2$, every hyperplane of $M_n(\mathbb K)$ contains an invertible matrix.
\end{theorem}
\begin{proof}
By Proposition~\ref{prop:trace} the hyperplane is $H=\{X:\operatorname{tr}(AX)=0\}$ for some $A\ne0$. Let $r=\operatorname{rank}A\ge1$.
By Gaussian elimination on rows and columns there are invertible $P,Q$ with $A=PJ_rQ$, where $J_r$ is the diagonal matrix
with $r$ ones followed by zeros. Let $C$ be the permutation matrix of the cycle $e_1\mapsto e_2\mapsto\dots\mapsto e_n\mapsto
e_1$; it is invertible, and since $n\ge2$ its diagonal is zero. Put $X=Q^{-1}CP^{-1}$, which is invertible. Then
$\operatorname{tr}(AX)=\operatorname{tr}(PJ_rCP^{-1})=\operatorname{tr}(J_rC)=\sum_{i\le r}c_{ii}=0$, so $X\in H$.
\end{proof}

For $n=1$ the statement fails, since the only hyperplane of $M_1(\mathbb K)=\mathbb K$ is $\{0\}$. The theorem says that
the singular matrices, although they form a set defined by one equation $\det X=0$, contain no hyperplane; Exercise
\ref{ex:singularspace} shows how large a subspace of singular matrices can be.

The third instance shows that the dimension argument can find points as well as bound families.

\begin{proposition}[unisolvent points]\label{prop:unisolvent}
Let $V$ be an $n$-dimensional space of functions $S\to\mathbb K$. There exist points $s_1,\dots,s_n\in S$ such that every
choice of values $y_1,\dots,y_n$ is taken at the $s_i$ by exactly one function of $V$.
\end{proposition}
\begin{proof}
The evaluations $\varepsilon_s(f)=f(s)$, $s\in S$, are linear forms on $V$. The subspace of $V^*$ they span has annihilator
$\{f\in V:f(s)=0\text{ for all }s\}=\{0\}$, so by Theorem~\ref{thm:annihilator}(a) it is all of $V^*$. Extract a basis
$\varepsilon_{s_1},\dots,\varepsilon_{s_n}$ of $V^*$ (Theorem~\ref{thm:dim}). Then the map $V\to\mathbb K^n$,
$f\mapsto(f(s_1),\dots,f(s_n))$, is injective (an $f$ killed by a basis of $V^*$ is killed by all of $V^*$, so is $0$),
hence bijective.
\end{proof}

For $V=\mathbb K[X]_{\le n}$, any $n+1$ distinct points work, and we recover Lagrange interpolation. For the trigonometric
polynomials $\sum_{|k|\le N}c_ke^{ikx}$ any $2N+1$ distinct points of $[0,2\pi)$ work (Exercise~\ref{ex:trig}). In two
variables things are subtler: six points on a circle are never unisolvent for the polynomials of degree at most $2$,
since the equation of the circle vanishes at all of them. The proposition guarantees that \emph{some} six points work.

\begin{beauty}
Grassmann's insight, that a space should be studied through its subspaces and the maps out of it, turns into a
symmetry of the whole theory. Every subspace $W$ of $V$ has a partner $W^\circ$ in $V^*$; sums become intersections,
inclusions are reversed, dimensions are complemented, and the partner of the partner is the original. Every theorem
about subspaces and maps has a dual theorem, proved for free. The same reflection, in projective geometry, exchanges
points and lines (Tome~S2, Chapter~4), and in optimisation, it exchanges a problem and its certificate of optimality.
\end{beauty}

\section{Exercises}

\begin{exercise}[$\star$]\label{ex:logs}
Show that $\log2$, $\log3$ and $\log5$ are linearly independent over $\Q$. More generally, show that the logarithms of
the primes are linearly independent over $\Q$, and deduce again that $\R$ is infinite-dimensional over $\Q$.
\end{exercise}

\begin{exercise}[$\star$]\label{ex:grassmann}
Let $U,W$ be subspaces of dimension $6$ in a space of dimension $10$. What are the possible dimensions of $U\cap W$?
Show that $U+W\cong(U\oplus W)/D$ for a subspace $D$ of $U\oplus W$ isomorphic to $U\cap W$, and deduce Grassmann's
formula from Theorem~\ref{thm:quotient}.
\end{exercise}

\begin{exercise}[$\star$]\label{ex:quotientdual}
Let $W\subseteq V$. Show that $V^*/W^\circ\cong W^*$ canonically (restrict forms to $W$).
\end{exercise}

\begin{exercise}[$\star\star$]\label{ex:traceform}
(a) Show that the linear forms $\varphi$ on $M_n(\mathbb K)$ with $\varphi(XY)=\varphi(YX)$ for all $X,Y$ are the multiples of the
trace. (b) Deduce that the span of the commutators $XY-YX$ is the hyperplane of matrices of trace $0$.
\end{exercise}

\begin{exercise}[$\star$]\label{ex:symanti}
For the pairing $(A,X)\mapsto\operatorname{tr}(AX)$ on $M_n(\R)$, find the annihilator of the symmetric matrices and of the
upper triangular matrices.
\end{exercise}

\begin{exercise}[$\star\star$]\label{ex:shift}
Find all linear forms $\varphi$ on $\R[X]_{\le n}$ such that $\varphi(P(X+1))=\varphi(P(X))$ for every $P$.
\end{exercise}

\begin{exercise}[$\star\star$]\label{ex:quadrature}
(a) Show that for distinct $x_0,\dots,x_n\in[0,1]$ there are unique weights $w_0,\dots,w_n$ such that
$\int_0^1P(t)\dd t=\sum_iw_iP(x_i)$ for all $P$ of degree at most $n$. (b) Compute them for $n=2$ and the points
$0,\frac12,1$: this is Simpson's rule. Show that it is in fact exact up to degree $3$.
\end{exercise}

\begin{exercise}[$\star$]\label{ex:trig}
Show that every $2N+1$ distinct points of $[0,2\pi)$ are unisolvent for the trigonometric polynomials
$\sum_{|k|\le N}c_ke^{ikx}$ with complex coefficients.
\end{exercise}

\begin{exercise}[$\star\star$]\label{ex:lights3}
Show without a computer that every configuration of the $3\times3$ Lights Out game can be switched off, by finding, for
each single light, a set of buttons that toggles exactly that light. Why does this suffice?
\end{exercise}

\begin{exercise}[$\star\star$]\label{ex:eventown}
(Reverse Oddtown.) Let $C_1,\dots,C_m\subseteq\{1,\dots,n\}$ have even sizes and pairwise odd intersections. Show that
$m\le n+1$, and that $m\le n$ when $m$ is even. (Hint: the Gram matrix over $\mathbb F_2$ is $J-I$.)
\end{exercise}

\begin{exercise}[$\star\star$]\label{ex:singularspace}
(a) Find a subspace of $M_n(\mathbb K)$ of dimension $n(n-1)$ consisting of singular matrices. (b) Show that a subspace
of $M_n(\mathbb K)$ in which every non-zero matrix is invertible has dimension at most $n$, and find one of dimension $2$
in $M_2(\R)$.
\end{exercise}

\begin{exercise}[$\star\star$]\label{ex:hyperplanes}
Let $\varphi,\varphi_1,\dots,\varphi_k\in V^*$. Show that $\varphi$ is a linear combination of $\varphi_1,\dots,\varphi_k$ if and only if
$\ker\varphi_1\cap\dots\cap\ker\varphi_k\subseteq\ker\varphi$.
\end{exercise}

\begin{exercise}[$\star\star\star$]\label{ex:hamel}
Assume, as Zorn's lemma allows (Tome~6, Chapter~1), that $\R$ has a basis $B$ over $\Q$. Show that there is a function
$f:\R\to\R$ with $f(x+y)=f(x)+f(y)$ for all $x,y$ that is not of the form $f(x)=cx$. Show that such an $f$ is unbounded
on every interval.
\end{exercise}

\begin{exercise}[$\star\star\star$]\label{ex:fisher}
(Fisher's inequality, the uniform case.) Let $C_1,\dots,C_m$ be distinct subsets of $\{1,\dots,n\}$ such that all pairwise
intersections have the same size $\lambda\ge1$. Show that $m\le n$. (Work over $\R$ with the Gram matrix of the indicator
vectors.)
\end{exercise}

\section{Hints}

\begin{hint}{ex:logs}
A rational relation can be cleared of denominators; exponentiate and use unique factorisation.
\end{hint}
\begin{hint}{ex:grassmann}
Consider the map $U\oplus W\to V$, $(u,w)\mapsto u+w$, and its kernel.
\end{hint}
\begin{hint}{ex:quotientdual}
The restriction map $V^*\to W^*$ is onto (extend forms using a complement); find its kernel.
\end{hint}
\begin{hint}{ex:traceform}
Write $\varphi=\operatorname{tr}(A\,\cdot)$ and test the condition on elementary matrices $E_{ij}$. For (b), compute the
annihilator of the span of the commutators.
\end{hint}
\begin{hint}{ex:symanti}
$\operatorname{tr}(AS)=0$ for all symmetric $S$: test on $E_{ii}$ and $E_{ij}+E_{ji}$.
\end{hint}
\begin{hint}{ex:shift}
Write the condition as $\varphi\circ(T-I)=0$ with $T(P)=P(X+1)$, so $\varphi\in(\operatorname{Im}(T-I))^\circ$.
\end{hint}
\begin{hint}{ex:quadrature}
$P\mapsto\int_0^1P$ is a linear form; expand it in the dual basis of the evaluations.
\end{hint}
\begin{hint}{ex:trig}
Multiply by $e^{iNx}$ and substitute $z=e^{ix}$.
\end{hint}
\begin{hint}{ex:lights3}
By symmetry three cases suffice: a corner, an edge and the centre. The solutions form a vector space.
\end{hint}
\begin{hint}{ex:eventown}
Compute $(J-I)^2$ over $\mathbb F_2$, and note that the rank of a Gram matrix of vectors in $\mathbb F_2^n$ is at most $n$.
\end{hint}
\begin{hint}{ex:singularspace}
(a) Matrices with a zero first column. (b) A subspace of dimension $n+1$ meets the subspace of matrices with zero first
column.
\end{hint}
\begin{hint}{ex:hyperplanes}
Consider $f:V\to\mathbb K^k$, $f(v)=(\varphi_1(v),\dots,\varphi_k(v))$ and the transpose of $f$.
\end{hint}
\begin{hint}{ex:hamel}
Define $f$ on the basis arbitrarily and extend $\Q$-linearly. For the second part, an additive function bounded on an
interval is linear.
\end{hint}
\begin{hint}{ex:fisher}
The Gram matrix is $\lambda J+D$ with $D$ diagonal, $d_i=|C_i|-\lambda\ge0$, and at most one $d_i$ is zero.
\end{hint}

\section{Solutions}

\begin{solution}{ex:logs}
\idea A linear relation between logarithms is a multiplicative relation between numbers, and unique factorisation
forbids those.
\writeup Suppose $\sum_i q_i\log p_i=0$ with rationals $q_i$ not all zero, $p_i$ distinct primes. Multiplying by a common
denominator we may assume $q_i\in\Z$. Moving the negative terms to the other side and exponentiating gives
$\prod_{q_i>0}p_i^{q_i}=\prod_{q_i<0}p_i^{-q_i}$, two factorisations of the same integer into distinct primes. By unique
factorisation all exponents are $0$, a contradiction. Since there are infinitely many primes, $\R$ contains an infinite
free family over $\Q$, so it is infinite-dimensional (a space of finite dimension $n$ contains no $n+1$ independent
vectors by Lemma~\ref{lem:steinitz}).
\end{solution}

\begin{solution}{ex:grassmann}
\idea Grassmann's formula is a rank--nullity statement for the addition map.
\writeup The map $s:U\oplus W\to V$, $s(u,w)=u+w$, has image $U+W$ and kernel $D=\{(x,-x):x\in U\cap W\}$, which is
isomorphic to $U\cap W$. By Theorem~\ref{thm:quotient}(c), $U+W\cong(U\oplus W)/D$, hence
$\dim(U+W)=\dim U+\dim W-\dim(U\cap W)$. With $\dim U=\dim W=6$ and $\dim(U+W)\le10$ we get $\dim(U\cap W)\ge2$, and of
course $\dim(U\cap W)\le6$. All values $2,\dots,6$ occur: with a basis $e_1,\dots,e_{10}$, take $U=\operatorname{span}(e_1,\dots,e_6)$
and $W=\operatorname{span}(e_{7-k},\dots,e_{12-k})$ for $k=2,\dots,6$, whose intersection is spanned by $e_{7-k},\dots,e_6$.
\end{solution}

\begin{solution}{ex:quotientdual}
\idea Restriction is a natural map $V^*\to W^*$; compute its kernel and image.
\writeup Let $r:V^*\to W^*$, $r(\varphi)=\varphi|_W$. Its kernel is $W^\circ$ by definition. It is onto: given $\psi\in W^*$,
choose a complement $C$ of $W$ and set $\varphi(w+c)=\psi(w)$. By the first isomorphism theorem, $V^*/W^\circ\cong W^*$. No
choice enters the isomorphism itself; the complement was only used to prove surjectivity.
\end{solution}

\begin{solution}{ex:traceform}
\idea The trace pairing turns a form into a matrix $A$; the condition becomes a condition on $A$ that elementary
matrices make concrete.
\writeup (a) By Proposition~\ref{prop:trace}, $\varphi(X)=\operatorname{tr}(AX)$. The condition says
$\operatorname{tr}(AXY)=\operatorname{tr}(AYX)=\operatorname{tr}(XAY)$, that is, $\operatorname{tr}((AX-XA)Y)=0$ for all $Y$, so
$AX=XA$ for all $X$ (Proposition~\ref{prop:trace} again). Taking $X=E_{ij}$: $AE_{ij}$ has column $j$ equal to column $i$ of
$A$, and $E_{ij}A$ has row $i$ equal to row $j$ of $A$, the rest zero. Comparing entries, $a_{ki}=0$ for $k\ne i$ and
$a_{ii}=a_{jj}$. So $A=cI$ and $\varphi=c\operatorname{tr}$. Conversely the trace satisfies $\operatorname{tr}(XY)=\operatorname{tr}(YX)$.
(b) Let $K$ be the span of the commutators. A form vanishes on $K$ exactly when $\varphi(XY)=\varphi(YX)$ for all $X,Y$, so
$K^\circ=\mathbb K\operatorname{tr}$ has dimension $1$. By Theorem~\ref{thm:annihilator}, $\dim K=n^2-1$ and $K=K^{\circ\circ}$
is the set of matrices killed by the trace.
\end{solution}

\begin{solution}{ex:symanti}
\idea Test the condition on a basis of the subspace.
\writeup $\operatorname{tr}(AE_{ii})=a_{ii}$ and $\operatorname{tr}(A(E_{ij}+E_{ji}))=a_{ji}+a_{ij}$. So $A$ annihilates the
symmetric matrices if and only if $a_{ii}=0$ and $a_{ij}=-a_{ji}$: the annihilator is the space of antisymmetric
matrices (dimension $n(n-1)/2=n^2-n(n+1)/2$, as it should be). For the upper triangular matrices,
$\operatorname{tr}(AE_{ij})=a_{ji}$ for $i\le j$ must vanish: the annihilator is the strictly upper triangular matrices.
\end{solution}

\begin{solution}{ex:shift}
\idea The invariance says that $\varphi$ kills every $P(X+1)-P(X)$; so $\varphi$ lies in the annihilator of the image of
the difference operator, and we only need that image.
\writeup Let $\Delta P=P(X+1)-P(X)$. The condition is $\varphi\circ\Delta=0$, that is, $\varphi\in(\operatorname{Im}\Delta)^\circ$.
$\Delta$ lowers degrees by exactly one ($\Delta X^k=kX^{k-1}+\dots$), so it maps $\R[X]_{\le n}$ into $\R[X]_{\le n-1}$, and
its kernel is the constants (a non-constant $P$ has $\deg\Delta P=\deg P-1\ge0$). By rank--nullity its image has dimension
$n$, so it is all of $\R[X]_{\le n-1}$. The annihilator of $\R[X]_{\le n-1}$ in the dual of $\R[X]_{\le n}$ is spanned by
the form ``coefficient of $X^n$''. The answer: $\varphi(P)=c\cdot[X^n]P$, $c\in\R$.
\end{solution}

\begin{solution}{ex:quadrature}
\idea Integration is a linear form; every linear form is a combination of the dual basis (here, evaluations), with
coefficients obtained by evaluating it on the basis (here, the Lagrange polynomials).
\writeup (a) The evaluations $\varepsilon_i$ form a basis of $(\R[X]_{\le n})^*$ (Example~\ref{ex:lagrange}), so the form
$I(P)=\int_0^1P$ is uniquely $\sum_iw_i\varepsilon_i$, and Proposition~\ref{prop:dualbasis} gives $w_i=I(L_i)$.
(b) With $x_0,x_1,x_2=0,\frac12,1$: $L_0=2(X-\frac12)(X-1)$, $L_1=-4X(X-1)$, $L_2=2X(X-\frac12)$, and integrating,
$w_0=\frac16$, $w_1=\frac23$, $w_2=\frac16$. So $\int_0^1P\approx\frac16\bigl(P(0)+4P(\frac12)+P(1)\bigr)$. For $P=X^3$:
$\frac16(0+\frac12+1)=\frac14=\int_0^1t^3\dd t$. Since both sides are linear and agree on $1,X,X^2,X^3$, the rule is exact
up to degree $3$ (the symmetry of the points around $\frac12$ is the reason; Chapter~4 explains Gauss's optimal choice).
\end{solution}

\begin{solution}{ex:trig}
\idea A trigonometric polynomial of degree $N$ is, after multiplying by $e^{iNx}$, an ordinary polynomial of degree $2N$
in $z=e^{ix}$.
\writeup The space has dimension $2N+1$ (the $e^{ikx}$, $|k|\le N$, are independent; see Chapter~9 or Tome~4,
Chapter~6). If $f=\sum c_ke^{ikx}$ vanishes at $2N+1$ distinct points $x_j$, then $Q(z)=\sum_kc_kz^{k+N}$, a polynomial of
degree at most $2N$, vanishes at the $2N+1$ distinct points $z_j=e^{ix_j}$ ($x\mapsto e^{ix}$ is injective on $[0,2\pi)$),
so $Q=0$ and $f=0$. The evaluation map to $\C^{2N+1}$ is thus injective, hence bijective.
\end{solution}

\begin{solution}{ex:lights3}
\idea The configurations that can be switched off are those of the form $Ax$: they form a subspace, the image of $A$. If
this subspace contains the nine single lights, it contains every configuration. Guessing button sets is fragile; the
reliable way is to solve the linear system, using the symmetry of the board to cut down the unknowns.
\writeup Number the cells $(r,c)$, $1\le r,c\le3$. To toggle the corner $(1,1)$ alone, look for a solution $x$ of
$Ax=e_{(1,1)}$ that is symmetric under the reflection in the main diagonal (this reflection maps the problem to itself, so
it maps solutions to solutions; we look for a symmetric one and will find it). Write $a=x_{11}$, $b=x_{12}=x_{21}$,
$c=x_{13}=x_{31}$, $d=x_{22}$, $e=x_{23}=x_{32}$, $f=x_{33}$. Cell $(r,c)$ is toggled by the presses at itself and its
neighbours, so the equations at $(1,1)$, $(1,2)$, $(1,3)$, $(2,2)$, $(2,3)$, $(3,3)$ read, in $\mathbb F_2$,
\[
\begin{gathered}a+2b=1,\qquad a+b+c+d=0,\qquad b+c+e=0,\\ 2b+d+2e=0,\qquad c+d+e+f=0,\qquad 2e+f=0 .\end{gathered}
\]
Since $2=0$: $a=1$, $d=0$, $f=0$; then $b+c=1$ and $b+c+e=0$ give $e=1$; $c+d+e+f=0$ gives $c=1$, hence $b=0$. So:
\emph{press $(1,1),(1,3),(3,1),(2,3),(3,2)$}. Check directly: $(1,1)$ is toggled only by its own press; each of the eight
other cells is toggled exactly twice (for instance $(2,2)$ by $(2,3)$ and $(3,2)$, and $(3,3)$ by $(2,3)$ and $(3,2)$).
The same method gives the other two cases: pressing $(2,2),(3,1),(3,2),(3,3)$ toggles only the edge cell $(1,2)$, and
pressing $(1,2),(2,1),(2,2),(2,3),(3,2)$ toggles only the centre, which is toggled five times while every outer cell is
toggled twice or not at all. The symmetries of the square carry these three solutions to solutions for all nine cells.
So $\operatorname{Im}A$ contains a basis of $\mathbb F_2^9$, $A$ is invertible, and every configuration can be switched
off, in exactly one way up to pressing a button twice.
\end{solution}

\begin{solution}{ex:eventown}
\idea Copy the proof of Oddtown: the Gram matrix of the indicator vectors over $\mathbb F_2$ is now $J-I$, and its rank
bounds $m$ from below in terms of $n$.
\writeup Let $v_i\in\mathbb F_2^n$ be the indicator vectors and $G=(v_i\cdot v_j)$ their Gram matrix, so $G=J-I$ ($J$ the
all-ones $m\times m$ matrix). If $M$ is the $m\times n$ matrix with rows $v_i$, then $G=MM^{\mathsf T}$, so
$\operatorname{rank}G\le\operatorname{rank}M\le n$. Over $\mathbb F_2$, $J^2=mJ$, so $(J-I)^2=J^2-2J+I=mJ+I$. If $m$ is
even this is $I$, so $G$ is invertible, $\operatorname{rank}G=m$ and $m\le n$. If $m$ is odd, the rows of $G$ sum to
$(m-1)(1,\dots,1)=0$, so $\operatorname{rank}G\le m-1$; on the other hand the top-left $(m-1)\times(m-1)$ block of $G$ is
$J_{m-1}-I_{m-1}$ with $m-1$ even, invertible by the first case, so $\operatorname{rank}G=m-1\le n$ and $m\le n+1$.
\end{solution}

\begin{solution}{ex:singularspace}
\idea For (a), force a common kernel vector. For (b), any large subspace must meet the space found in (a), by
Grassmann's formula.
\writeup (a) The matrices whose first column is zero form a subspace $Z$ of dimension $n^2-n$, and they all kill $e_1$,
so they are singular. (b) If $L$ is a subspace with $\dim L\ge n+1$ then $\dim(L\cap Z)\ge(n+1)+(n^2-n)-n^2=1$, so $L$
contains a non-zero singular matrix. In $M_2(\R)$, $L=\left\{\begin{pmatrix}a&-b\\b&a\end{pmatrix}\right\}$ works:
$\det=a^2+b^2>0$ unless $a=b=0$. (This copy of $\C$ inside $M_2(\R)$ is the reason; in $M_4(\R)$ the quaternions give a
$4$-dimensional example, Tome~2, Chapter~12.)
\end{solution}

\begin{solution}{ex:hyperplanes}
\idea ``$\varphi$ is a combination of the $\varphi_i$'' says $\varphi\in\operatorname{Im}f^{\mathsf T}$ for the map $f$ that
records all the $\varphi_i$; Theorem~\ref{thm:transpose}(b) identifies that image with an annihilator.
\writeup Let $f:V\to\mathbb K^k$, $f(v)=(\varphi_i(v))_i$. In the dual basis $(\varepsilon_i)$ of $\mathbb K^k$,
$f^{\mathsf T}(\varepsilon_i)=\varepsilon_i\circ f=\varphi_i$, so $\operatorname{Im}f^{\mathsf T}=\operatorname{span}(\varphi_1,\dots,\varphi_k)$.
By Theorem~\ref{thm:transpose}(b) this is $(\ker f)^\circ=(\ker\varphi_1\cap\dots\cap\ker\varphi_k)^\circ$. So $\varphi$ is a
combination of the $\varphi_i$ if and only if it vanishes on $\bigcap\ker\varphi_i$.
\end{solution}

\begin{solution}{ex:hamel}
\idea An additive function is $\Q$-linear, and a $\Q$-linear map can be prescribed arbitrarily on a $\Q$-basis. Pick
values that break proportionality.
\writeup An additive $f$ satisfies $f(nx)=nf(x)$ for $n\in\Z$ (induction), hence $f(qx)=qf(x)$ for $q\in\Q$ (write
$q=a/b$ and use $bf(qx)=f(ax)=af(x)$): additive means $\Q$-linear. Choose two distinct elements $b_1,b_2$ of the basis $B$
and define $f(b_1)=b_1$, $f(b)=0$ for $b\ne b_1$, extended $\Q$-linearly. Then $f$ is additive; if $f(x)=cx$, then
$c=f(b_1)/b_1=1$ and $c=f(b_2)/b_2=0$, a contradiction.

Now suppose $|f|\le M$ on an interval $[u,u+\delta]$. For $|t|\le\delta$, write $t=s-s'$ with $s,s'\in[u,u+\delta]$
(take $s=u+\max(t,0)$, $s'=u+\max(-t,0)$), so $|f(t)|\le2M$. Let $c=f(1)$ and $g(x)=f(x)-cx$, additive, bounded by
$2M+|c|\delta$ on $[-\delta,\delta]$, and zero on $\Q$. For any $x$ and rational $r$ with $|x-r|\le\delta/n$:
$|g(x)|=|g(x-r)|=\frac1n|g(n(x-r))|\le\frac{2M+|c|\delta}n$. Letting $n\to\infty$, $g=0$ and $f(x)=cx$. So a non-linear
additive function is unbounded on every interval; its graph is in fact dense in the plane.
\end{solution}

\begin{solution}{ex:fisher}
\idea The Oddtown trick over $\R$: compute the Gram matrix of the indicator vectors, and show it is invertible; then
$m\le\operatorname{rank}\le n$.
\writeup Let $v_i\in\R^n$ be the indicator vectors and $G=(v_i\cdot v_j)$. Then $G=\lambda J+D$, where $D$ is diagonal with
$d_i=|C_i|-\lambda$. Since $C_i\supseteq C_i\cap C_j$, $d_i\ge0$; and if $d_i=d_j=0$ for $i\ne j$ then $C_i=C_i\cap C_j=C_j$,
contradicting distinctness. So at most one $d_i$ vanishes. For $x\in\R^m$,
\[
x^{\mathsf T}Gx=\lambda\Bigl(\sum_ix_i\Bigr)^2+\sum_id_ix_i^2 .
\]
If this is $0$, then $\sum_id_ix_i^2=0$ forces $x_i=0$ whenever $d_i>0$, so at most one coordinate, say $x_k$, can be
non-zero, and then $\lambda x_k^2=0$ forces $x_k=0$. Hence $x^{\mathsf T}Gx>0$ for $x\ne0$, $G$ is invertible, and since
$G=MM^{\mathsf T}$ with $M$ of size $m\times n$, $m=\operatorname{rank}G\le n$. (In design theory this is Fisher's inequality of 1940:
a block design has at least as many blocks as points. Tome~C1, Chapter~11 proves the non-uniform version and much more.)
\end{solution}

\section*{Sources and further reading}

\begin{itemize}
\item H.~Grassmann, \emph{Die lineale Ausdehnungslehre}, Wigand, Leipzig, 1844; English translation by L.~C.~Kannenberg,
\emph{A New Branch of Mathematics}, Open Court, 1995. The story of its reception is told in D.~Fearnley-Sander,
\emph{Hermann Grassmann and the creation of linear algebra}, American Mathematical Monthly 86 (1979), 809--817.
\item G.~Peano, \emph{Calcolo geometrico secondo l'Ausdehnungslehre di H.~Grassmann}, Bocca, Turin, 1888. Chapter~IX gives
the axioms of a vector space.
\item S.~Axler, \emph{Linear Algebra Done Right}, 4th ed., Springer, 2024. Chapter~3 for quotients, duality and the
transpose, in the basis-free spirit of this chapter.
\item P.~R.~Halmos, \emph{Finite-Dimensional Vector Spaces}, 2nd ed., Van Nostrand, 1958. The classic treatment of
annihilators and the double dual.
\item L.~Babai and P.~Frankl, \emph{Linear Algebra Methods in Combinatorics}, preliminary version, University of Chicago,
1992. Chapter~1 opens with Oddtown and Eventown.
\item E.~R.~Berlekamp, \emph{On subsets with intersections of even cardinality}, Canadian Mathematical Bulletin 12 (1969),
471--474.
\item K.~Sutner, \emph{Linear cellular automata and the Garden-of-Eden}, Mathematical Intelligencer 11 (1989), 49--53, for
the linear algebra of Lights Out on general boards.
\item H.~Flanders, \emph{On spaces of linear transformations with bounded rank}, Journal of the London Mathematical Society
37 (1962), 10--16: a subspace of singular $n\times n$ matrices has dimension at most $n(n-1)$, which shows that
Exercise~\ref{ex:singularspace}(a) is optimal.
\item R.~A.~Fisher, \emph{An examination of the different possible solutions of a problem in incomplete blocks}, Annals of
Eugenics 10 (1940), 52--75.
\end{itemize}
